% !TeX root = ../../Book3.tex
% 纲要标题：## 第3章　Noether 性、Artin 性与 Nakayama 引理
% 纲要位置：工作记录/计划纲要.md，第 1977 行。

\chapter{Noether 性、Artin 性与 Nakayama 引理}
\label{b3:ch:03}
\BookChapterNumber{b3:ch:03}

\begin{chapterabstract}
有限生成与有限展示使模能够用有限数据描述，Noether 条件进一步保证子模仍有这样的描述。本章以行列式技巧证明交换环上的 Nakayama 引理，进而用剩余向量空间计算局部模的最少生成元数。对交换 Artin 环，则先证明根的幂零性与有限长度，再得到局部直积分解和零维 Noether 环的刻画。最后比较商、局部化、有限代数与子环中的有限性行为。
\end{chapterabstract}

\begin{learninggoals}
\begin{itemize}
\item 区分有限生成、有限展示、Noether 性、Artin 性与有限长度。
\item 能在具体证明中指出行列式技巧和 Nakayama 引理的假设。
\item 用 \(M/\mathfrak mM\) 选择极小生成组，并避免把生成组误认为自由基。
\item 分解交换 Artin 环，并用反例说明零维刻画为何需要 Noether 条件。
\end{itemize}
\end{learninggoals}

固定素数 \(p\)，设 \(R=\mathbb Z_{(p)}\)，把有理数域 \(\mathbb Q\) 当作 \(R\) 模。它非零，却满足 \(p\mathbb Q=\mathbb Q\)，所以 \(\mathbb Q/p\mathbb Q=0\)。因此“模掉极大理想后为零”不能单独推出原模为零；有限生成是 Nakayama 引理不可少的条件。反过来，对一个有限生成局部模，剩余域上的向量空间又准确提示最少需要多少生成元。

\ProseParagraphBreak

有限性还会改变理想与子模的行为。若一切理想都有限生成，任意有限生成模的子模也能继续由有限数据控制；若升链必定停止，反复取理想或子模就不会无限产生新的条件。对 Artin 环，还要把这种停止性和局部商的结构合起来，才能说明它如何分解成有限个局部部分。

\ProseParagraphBreak

预备知识为前两章的理想、素谱、局部化与局部环，以及第二卷第12.2节的有限长度与 Jordan–Hölder 定理\footnote{若尔当（Marie Ennemond Camille Jordan，1838-1922），法国数学家，研究群论与 Galois 理论。赫尔德（Ludwig Otto Hölder，1859-1937），德国数学家，研究群的合成列与实数的公理化。}。Hilbert 基定理的主要证明已在第一卷定理12.4.10给出，本章据此发展模论后果。

\ProseParagraphBreak
Matsumura 的《Commutative Ring Theory》\cite[Theorems 2.1--2.3, pp. 7--9; Theorem 3.2, p. 16]{Matsumura1989CommutativeRingTheory}可用于对照行列式技巧、局部生成元和交换 Artin 环的证明。

% !TeX root = ../../Book3.tex
% 纲要标题：### 3.1 有限性条件
% 纲要位置：工作记录/计划纲要.md，第 1979 行。
% 纲要细目（写作范围）：
% * 有限生成与有限展示；基本关系生成全部关系及有限展示对生成组的独立性
% * Noether 环上的有限生成模；升链条件与子模有限生成的等价性
% * Hilbert 基定理的模论后果

\section{有限性条件}
\label{b3:sec:0301}

有限个生成元给出模中元素的表达方式，却不保证生成元之间的关系能由有限组关系生成。再要求每个子模都有限生成，就得到 Noether 条件。这些条件控制不同的有限性，不能混用；Hilbert 基定理则把底环的 Noether 性传给有限变量多项式环，使多项式环上的有限生成模及其子模也具有有限描述。

\subsection{有限生成与有限展示}
\label{b3:subsec:0301:01}

\begin{definition}\label{b3:def:finite-presentation}
设 \(R\) 为交换含幺环，\(M\) 为 \(R\) 模。若存在非负整数 \(b\) 及元素 \(m_1,\ldots,m_b\in M\)，使每个 \(m\in M\) 都能写成
\(\sum_{i=1}^b a_i m_i\)，其中 \(a_i\in R\)，则称 \(M\) 为\term[youxian-shengcheng-mo]{有限生成模}[finitely generated module]。若存在非负整数 \(a,b\) 及 \(R\) 模正合列
\[
R^a\xrightarrow{A}R^b\longrightarrow M\longrightarrow0,
\]
则称 \(M\) 为\term[youxian-zhanshi-mo]{有限展示模}[finitely presented module]。该正合列给出一个有限展示：满射把 \(R^b\) 的标准基送到 \(M\) 的生成元，其核由矩阵 \(A\) 的列生成，称为这组生成元的关系模。
\end{definition}
有限展示模也称有限呈示模\termidx[youxian-chengshi-mo]{有限呈示模}；本书以“展示”为主称。“有限生成”限制的是生成元数目，而非底集合的大小：例如 \(\mathbb Z\) 作为整数模由 \(1\) 生成，任意非零有限维实向量空间也有限生成，二者的底集合却都无限。

\ProseParagraphBreak
若 \(m_1,\ldots,m_b\) 生成 \(M\)，自由模 \(R^b\) 的标准基 \(e_i\) 可以分别送到 \(m_i\)，得到满射 \(p:R^b\to M\)。向量 \(\sum_i r_i e_i\) 属于 \(\ker p\)，恰好表示关系 \(\sum_i r_i m_i=0\)，所以这个核记录全部关系。例如，对 \(f,g\in R\)，有有限展示
\[
R^2\xrightarrow{(f\ g)}R\longrightarrow R/(f,g)\longrightarrow0.
\]
其第一个映射把 \((a,b)\) 送到 \(af+bg\)，第二个映射为自然商映射。模 \(R/(f,g)\) 由 \(e=1+(f,g)\) 生成，\(fe=0\)、\(ge=0\) 是两条基本关系；任一关系 \(re=0\) 都由 \(r\in(f,g)\) 得到，因而是这两条基本关系的 \(R\) 线性组合。这里的有限性限制基本关系的数目，不要求全部关系只有有限条。

\ProseParagraphBreak
换一组有限生成元后，需要考虑相应的自由模满射及其核。下面的命题说明，一旦某组有限生成元的关系模有限生成，其他有限生成组也有这一性质；因此有限展示性不依赖生成组的选择。

\begin{proposition}\label{b3:prop:finite-presentation-kernel}
设 \(R\) 为交换含幺环，\(M\) 为 \(R\) 模。\(M\) 有限展示，当且仅当它有限生成且存在满射 \(R^b\rightarrow M\)，其中 \(b\ge0\)，其核有限生成。此时，对每个非负整数 \(n\)，满射 \(R^n\rightarrow M\) 的核也都有限生成。
\end{proposition}
\begin{proof}
第一句只是把关系矩阵的像写成核；反过来，对核选有限个生成元即可构成有限展示。证明最后一句时，设 \(p:R^b\rightarrow M\) 的核 \(L\) 有限生成，另有满射 \(q:R^n\rightarrow M\)。对 \(R^b\) 的每个标准基向量 \(e_i\)，由 \(q\) 满射选取 \(y_i\in R^n\)，使 \(q(y_i)=p(e_i)\)；规定 \(u(e_i)=y_i\)，便得到 \(u:R^b\rightarrow R^n\) 且 \(qu=p\)。反过来，对 \(R^n\) 的标准基向量 \(f_j\)，选取 \(z_j\in R^b\)，使 \(p(z_j)=q(f_j)\)；规定 \(v(f_j)=z_j\)，得到 \(v:R^n\rightarrow R^b\) 且 \(pv=q\)。这两个映射分别把一组生成元用另一组来表达。若 \(K=\ker q\)，则 \(u(L)\subseteq K\)，且
\(q(1-uv)=q-pv=0\)。对 \(x\in K\)，有 \(vx\in L\)，所以
\[
x=u(vx)+(1-uv)x.
\]
因此 \(K=u(L)+(1-uv)(R^n)\)，由 \(L\) 的有限生成元的像和有限个 \((1-uv)f_j\) 生成。
\end{proof}

循环模 \(R/I\) 总是有限生成；由这个命题，它有限展示当且仅当理想 \(I\) 有限生成。取
域 \(k\)，令 \(R=k[x_1,x_2,\ldots]\)、\(I=(x_1,x_2,\ldots)\)，便得到有限生成而非有限展示的模。若 \(I\) 由有限个多项式生成，把其中出现的变量集中在前 \(n\) 个内，再将这 \(n\) 个变量取零；所有生成元的像为零，\(x_{n+1}\) 的像却非零，矛盾。

\subsection{Noether 环上的有限生成模}
\label{b3:subsec:0301:02}

\begin{definition}\label{b3:def:noether-artin-module}
设 \(R\) 为交换含幺环，\(M\) 为 \(R\) 模。若 \(M\) 的每条子模升链 \(N_1\subseteq N_2\subseteq\cdots\) 最终稳定，则称 \(M\) 为 \term[Noether-mo]{Noether 模}[Noetherian module]。若每条子模降链最终稳定，则称 \(M\) 为 \term[Artin-mo]{Artin 模}[Artinian module]。若 \(R\) 作为自身的模满足相应条件，则分别称 \(R\) 为 \term[Noether-huan]{Noether 环}[Noetherian ring]与 \term[Artin-huan]{Artin 环}[Artinian ring]。
\end{definition}

一个模是 Noether 模，当且仅当它的每个子模有限生成。若某个子模不有限生成，逐次添入不能由已有元素生成的元素，就得到严格升链。反过来，若每个子模有限生成，对一条升链 \(N_1\subseteq N_2\subseteq\cdots\)，其并 \(N\) 是子模。若 \(N=0\)，链各项已全为零；否则选有限组生成元 \(m_1,\ldots,m_r\)，分别取 \(i_j\) 使 \(m_j\in N_{i_j}\)。令 \(n=\max\{i_1,\ldots,i_r\}\)，链的嵌套关系保证全部生成元都在 \(N_n\) 中，所以 \(N_n=N\)，以后各项也只能等于 \(N\)。

\begin{proposition}\label{b3:prop:noether-finite-modules}
设 \(R\) 为交换含幺环，\(0\rightarrow L\rightarrow M\rightarrow N\rightarrow0\) 为 \(R\) 模的短正合列。中间项 \(M\) 为 Noether 模，当且仅当两端 \(L,N\) 均为 Noether 模；Artin 模也满足同样结论。特别地，若 \(R\) 为 Noether 环，则每个有限生成 \(R\) 模及其子模都为 Noether 模，并且都有限展示。
\end{proposition}
\begin{proof}
子模的链也是原模的链，商模的链可以取逆像，所以中间项的链条件传给两端。反向设两端具有升链条件。对 \(M\) 的升链 \((M_i)\)，交链 \(M_i\cap L\) 及其在 \(N\) 中的像链都稳定。若 \(M_i\subseteq M_j\)，且交与像相同，对 \(x\in M_j\) 可取 \(y\in M_i\) 有相同的像，则 \(x-y\in M_j\cap L=M_i\cap L\)，所以 \(x\in M_i\)。因此原链稳定。降链在共同的稳定起点后，对反向包含使用同一个论证，即得 Artin 情形。

若 \(R\) 为 Noether 环，对 \(n\ge1\)，在向量末尾补一个零给出 \(R^{n-1}\to R^n\) 的单射，取末坐标则给出 \(R^n\to R\) 的满射，其核恰为前一个映射的像。因此有短正合列 \(0\rightarrow R^{n-1}\rightarrow R^n\rightarrow R\rightarrow0\)，从零模 \(R^0\) 开始归纳得到有限自由模为 Noether 模；其商即任意有限生成模仍为 Noether 模，所有子模也如此。对这些模选有限自由满射，核作为有限自由 Noether 模的子模有限生成，再用\cref{b3:prop:finite-presentation-kernel}得到有限展示。
\end{proof}

这个命题把第二卷第12.1节的链条件论证落实到交换环上。它还解释了\cref{b3:thm:hom-localization}在 Noether 环上为何只需假设源模有限生成：有限展示此时自动成立。若底环不满足 Noether 条件，就不能省略那项假设。

\subsection{Hilbert 基定理的模论后果}
\label{b3:subsec:0301:03}

第一卷定理12.4.10，即 Hilbert 基定理，已经证明：若 \(R\) 为 Noether 环，则有限变量多项式环 \(R[x_1,\ldots,x_n]\) 仍为 Noether 环。这里不重复其主证明，而将它与\cref{b3:prop:noether-finite-modules}合用。

\begin{corollary}\label{b3:cor:hilbert-module-consequences}
设 \(R\) 为交换含幺 Noether 环，\(A=R[a_1,\ldots,a_n]\) 为有限型交换 \(R\) 代数，其中 \(n\ge0\)。则 \(A\) 为 Noether 环。每个有限生成 \(A\) 模有有限展示，每个子模有限生成；有限矩阵给出的核、像及余核也都是有限生成 \(A\) 模。
\end{corollary}
\begin{proof}
代入同态 \(R[x_1,\ldots,x_n]\rightarrow A\)、\(x_i\mapsto a_i\) 为满射。由 Hilbert 基定理，来源是 Noether 环；商环的理想链取逆像，故商环 \(A\) 仍为 Noether 环。现在对 \(A\) 使用\cref{b3:prop:noether-finite-modules}，得到关于有限生成模和子模的断言。矩阵的核是有限自由模的子模，像与余核分别是有限自由模的同态像及商，也满足这些结论。
\end{proof}

例如，取域 \(k\)，在 \(A=k[x,y]/(y^2-x^3)\) 上，理想、理想之间的关系以及有限矩阵的解模，都能用有限数据表示。这个结论保证有限描述存在，却没有给出具体的生成元算法；当系数域上可以作精确计算时，第一卷附录F的 Gröbner 基方法和本卷\chapref{b3:ch:08}将进一步处理算法问题。

\ProseParagraphBreak
必须分清有限型代数与有限生成模。对域 \(k\)，多项式环 \(k[x]\) 由一个元素作为 \(k\) 代数生成，却有无限多个线性无关的单项式，因而不是有限生成的 \(k\) 模。上面的推论说有限生成的 \(A\) 模是 Noether 模，没有说它作为 \(R\) 模也有限生成。

\ProseParagraphBreak
模的有限长度控制另一种有限性：它要求存在有限合成列
\(0=M_0\subsetneq M_1\subsetneq\cdots\subsetneq M_\ell=M\)，其中每个相邻商 \(M_i/M_{i-1}\) 都是简单模；零模对应 \(\ell=0\) 的空合成列。第二卷定理12.2.4证明，有限长度等价于同时满足 Noether 与 Artin 条件。\cref{b3:tab:module-finiteness}将这些条件放在一起比较。

\begin{table}[htbp]
\centering
\caption[模的有限性条件]{模的有限性条件。设 \(R\) 为交换含幺环，\(M\) 为 \(R\) 模；除最后一行外，不附加底环的链条件。}
\label{b3:tab:module-finiteness}
\begin{tabularx}{\textwidth}{@{}l>{\raggedright\arraybackslash}X>{\raggedright\arraybackslash}p{0.29\textwidth}@{}}
\toprule
条件 & 所要求的有限性 & 与其他条件的关系 \\
\midrule
有限生成 & 有有限生成集 & 不一定有限展示或为 Noether 模 \\
有限展示 & 有有限生成集，全部关系均为有限组基本关系的 \(R\) 线性组合 & 推出有限生成；不一定为 Noether 模 \\
Noether & 每条子模升链最终稳定；等价地，每个子模都有限生成 & 推出有限生成；不一定为 Artin 模 \\
Artin & 每条子模降链最终稳定 & 不一定有限生成 \\
有限长度 & 有有限合成列 & 等价于 Noether 且 Artin；因而有限生成 \\
\midrule
\(R\) 为 Noether 环 & 对任意 \(R\) 模 \(M\)，有限生成、有限展示与 Noether 条件等价 & 不要求 \(M\) 为 Artin 模 \\
\bottomrule
\end{tabularx}
\end{table}

有限长度在一般底环上也不保证有限展示。上面的无限变量例子中，\(R/I\cong k\) 是简单 \(R\) 模，长度为 \(1\)，但理想 \(I\) 不有限生成，故 \(R/I\) 没有有限展示。

% !TeX root = ../../Book3.tex
% 纲要标题：### 3.2 Nakayama 引理
% 纲要位置：工作记录/计划纲要.md，第 1985 行。
% 纲要细目（写作范围）：
% * 行列式技巧；生成元组上的矩阵关系及标量作用的像
% * Jacobson 根版本与局部环版本
% * 极小生成元及剩余向量空间；区别极小、最少与自由基

\section{Nakayama 引理}
\label{b3:sec:0302}

在局部环上，将有限生成模按极大理想取商，就得到剩余域上的有限维向量空间。Nakayama 引理说明，向量空间中的生成元可以提升为原模的生成元。这里的有限性十分关键；证明它的行列式技巧，还将在整依赖中给出首一方程。

\subsection{行列式技巧}
\label{b3:subsec:0302:01}

考虑有限生成 \(R\) 模 \(M\) 满足 \(M=IM\) 的情形，其中 \(R\) 为交换含幺环、\(I\) 为理想。选取生成元 \(m_1,\ldots,m_n\)，每个生成元都能写成 \(m_i=\sum_j a_{ij}m_j\)、\(a_{ij}\in I\)。把这些关系合在一起，就有
\[
(I_n-A)\begin{pmatrix}m_1\\ \vdots\\m_n\end{pmatrix}=0,
\qquad A=(a_{ij}).
\]
于是可以利用伴随矩阵恒等式，寻找同时零化全部生成元的一个环元素，而不再分别处理这些关系。对于满足 \(u(M)\subseteq IM\) 的一般自同态，同一个方法也能产生零化多项式。

\begin{theorem}[行列式技巧]\label{b3:thm:determinant-trick}
设 \(R\) 为交换含幺环，\(M\) 为由 \(n\ge1\) 个元素生成的 \(R\) 模，\(I\subseteq R\) 为理想，
\(u:M\rightarrow M\) 为 \(R\) 线性自同态，并且 \(u(M)\subseteq IM\)。则存在
\[
P(T)=T^n+c_1T^{n-1}+\cdots+c_n,\qquad c_i\in I^i,
\]
使 \(P(u)=0\) 作为 \(M\) 的自同态成立。
\end{theorem}
\begin{proof}
选生成元 \(m_1,\ldots,m_n\)，写
\(u(m_i)=\sum_j a_{ij}m_j\)、\(a_{ij}\in I\)。令 \(A=(a_{ij})\in\operatorname{Mat}_n(R)\)。在自同态环 \(\operatorname{End}_R(M)\) 中，\(R\) 的标量作用像与 \(u\) 生成一个交换子环，简记为 \(R[u]\)。这里使用的是映射 \(R\to\operatorname{End}_R(M)\) 的像，不要求它单射，也不要求整个自同态环交换。将每个 \(a_{ij}\) 理解为对应的标量自同态，上述等式写成
\[
(uI_n-A)\begin{pmatrix}m_1\\ \vdots\\m_n\end{pmatrix}=0.
\]
这是 \(M^n\) 中的等式，列中的 \(m_i\) 是 \(M\) 的元素，并不是相对于一组自由基的坐标。矩阵的每个自同态作用在相应分量上，因此可以继续左乘矩阵。对 \(B=uI_n-A\)，在交换环 \(R[u]\) 上使用 \(\operatorname{adj}(B)B=\det(B)I_n\)，得到
\[
\det(uI_n-A)m_i=0\qquad(1\le i\le n).
\]
令 \(P(T)=\det(TI_n-A)\in R[T]\)。行列式的展开中，\(T^{n-i}\) 的系数每项都含 \(i\) 个 \(A\) 的元素，故 \(c_i\in I^i\)。由于 \(P(u)\) 零化全部生成元，它零化整个 \(M\)。
\end{proof}

这称为\term[hanglieshi-jiqiao]{行列式技巧}[determinant trick]。矩阵 \(A\) 未必唯一，所产生的首一多项式也不必唯一，但它确实零化给定自同态。若 \(M\) 自由且所取生成元是一组基，便得到熟悉的 Cayley–Hamilton 恒等式\footnote{凯莱（Arthur Cayley，1821-1895），英国数学家，研究矩阵、不变量和代数几何。哈密顿（William Rowan Hamilton，1805-1865），爱尔兰数学家及物理学家，创立四元数理论。}。上述一般版本可对照 Matsumura 的《Commutative Ring Theory》\cite[Theorem 2.1, pp. 7--8]{Matsumura1989CommutativeRingTheory}。

\ProseParagraphBreak
若 \(M=IM\)，对 \(u=\operatorname{id}_M\) 使用该定理，得到
\[
(1+c_1+\cdots+c_n)M=0.
\]
由于 \(c_i\in I^i\subseteq I\) 对 \(i\ge1\) 成立，令 \(a\coloneqq-(c_1+\cdots+c_n)\)，便有 \(a\in I\) 且 \((1-a)M=0\)。有限生成将逐个生成元的关系合成了一个统一的零化元。

\subsection{Jacobson 根版本与局部环版本}
\label{b3:subsec:0302:02}

\begin{definition}\label{b3:def:jacobson-radical}
设 \(R\) 为交换含幺环。\(R\) 的\term[Jacobson-gen]{Jacobson 根}[Jacobson radical]为
\[
J(R)=\bigcap_{\mathfrak m\;\text{极大}}\mathfrak m.
\]
零环采用空交为整个环的约定。
\end{definition}
\symidx{\(J(R)\)：交换环的 Jacobson 根}

若 \(a\in J(R)\)，则对每个 \(r\in R\) 都有 \(ra\in J(R)\)，所以 \(1-ra\) 不属于任何极大理想，因而是单位。反过来，若 \(a\notin J(R)\)，由交的定义可取极大理想 \(\mathfrak m\) 使 \(a\notin\mathfrak m\)。于是 \(\bar a\) 在域 \(R/\mathfrak m\) 中非零，取其逆并提升为某个 \(r\in R\)，便有 \(1-ra\in\mathfrak m\)，所以 \(1-ra\) 不是单位。因此 Jacobson 根的单位判据为
\[
a\in J(R)\quad\Longleftrightarrow\quad
1-ra\in R^\times\;\text{对每个}\;r\in R\;\text{成立}.
\]
在局部环上，\(J(R)\) 就是唯一的极大理想。

\begin{theorem}[Nakayama 引理]\label{b3:thm:nakayama}
设 \(R\) 为交换含幺环，\(M\) 为有限生成 \(R\) 模，\(I\) 为 \(R\) 中满足 \(I\subseteq J(R)\) 的理想，\(N\subseteq M\) 为子模。
\begin{enumerate}
\item 若 \(IM=M\)，则 \(M=0\)。
\item 若 \(N\subseteq M\) 且 \(M=N+IM\)，则 \(M=N\)。
\end{enumerate}
第二项实际上只需 \(M/N\) 有限生成，而不必要求 \(M\) 本身有限生成。特别地，在局部环 \((R,\mathfrak m,\kappa)\) 上，有限生成模满足
\(M/\mathfrak mM=0\) 当且仅当 \(M=0\)。
\end{theorem}
\begin{proof}
第一项中，若 \(M\ne0\)，选有限个生成元。行列式技巧给出 \(a\in I\) 使 \((1-a)M=0\)。由 \(I\subseteq J(R)\)，\(1-a\) 为单位，乘其逆得到 \(M=0\)，矛盾；零模情形本来成立。第二项对有限生成商模 \(Q=M/N\) 使用第一项，因为原等式恰好表示 \(IQ=Q\)。局部版本取 \(I=\mathfrak m\)。
\end{proof}

\term[Nakayama-yinli]{Nakayama 引理}[Nakayama's lemma]的这一交换环证明，也可见\cite[Theorem 2.2 and Corollary, p. 8]{Matsumura1989CommutativeRingTheory}。第二卷定理14.2.1则用 Jacobson 根的单位判据直接处理非交换模，没有把交换行列式用于非交换系数。

\ProseParagraphBreak
第二项中的 \(M=N+IM\) 恰好表示自然映射 \(N\to M/IM\) 满射。因此，当 \(M\) 有限生成且 \(I\subseteq J(R)\) 时，在 \(M/IM\) 中选出一组生成元，再任取它们到 \(M\) 的提升，这些提升便生成 \(M\)。在局部环上取 \(I=\mathfrak m\)，只需在剩余向量空间中选择生成元，就能得到原模的生成元。

\ProseParagraphBreak
两项假设都有明确作用。取素数 \(p\)，对局部环 \(R=\mathbb Z_{(p)}\) 的模 \(M=\mathbb Q\)，有 \(pM=M\ne0\)，因为 \(M\) 不有限生成：有限个有理数只能给出有统一上界的负 \(p\) 指数。若省去 \(I\subseteq J(R)\)，则 \(R=\mathbb Z\)、\(I=(2)\)、\(M=\mathbb Z/3\mathbb Z\) 满足 \(IM=M\ne0\)。有限生成与 Jacobson 根条件不能相互替代。

\begin{corollary}\label{b3:cor:finite-surjective-endomorphism}
设 \(R\) 为交换含幺环，\(M\) 为有限生成 \(R\) 模。则 \(M\) 的每个满 \(R\) 线性自同态都是同构。
\end{corollary}
\begin{proof}
设 \(u:M\rightarrow M\) 满。满射性是 \(u(M)=M\)；把 \(u\) 编码成一个标量的作用，就能将它改写成前面使用的 \(IM=M\)。为此赋予 \(M\) 一个 \(R[T]\) 模结构，令 \(T\) 作用为 \(u\)；原来的有限个 \(R\) 生成元也生成此 \(R[T]\) 模，并且 \((T)M=u(M)=M\)。对交换环 \(R[T]\) 及其理想 \((T)\) 使用上面的统一零化元结论，得到
\(\bigl(1-T\cdot g(T)\bigr)M=0\)，其中 \(g(T)\in R[T]\)。所以
\(u\,g(u)=g(u)\,u=\operatorname{id}_M\)，给出了双侧逆。
\end{proof}

\subsection{极小生成元与剩余向量空间}
\label{b3:subsec:0302:03}

一个生成组不可再删去元素，只说明其中每个元素对这组生成元都不可缺少；元素个数最少则需要与所有生成组比较。生成组是否为自由基，还要检查是否存在系数不全为零的线性关系。在局部环上的有限生成模中，前两个条件由剩余向量空间联系起来。

\begin{corollary}\label{b3:cor:local-generators}
设 \((R,\mathfrak m,\kappa)\) 为交换局部环，其中 \(\kappa=R/\mathfrak m\)，\(M\) 为有限生成 \(R\) 模，\(m_1,\ldots,m_r\in M\)。这些元素生成 \(M\)，当且仅当其剩余类生成
\(\kappa\) 向量空间 \(M/\mathfrak mM\)。它们构成不可再删去元素的极小生成组，当且仅当剩余类构成该向量空间的一组基。因此所有极小生成组的元素个数都等于
\[
\mu_R(M)\coloneqq\dim_\kappa(M/\mathfrak mM).
\]
\end{corollary}
\begin{proof}
正向由取商得到。反向令 \(N=\sum_i Rm_i\)，剩余类生成意味着 \(M=N+\mathfrak mM\)，由\cref{b3:thm:nakayama}有 \(M=N\)。若剩余类为基，删去任一项便不再生成剩余向量空间，故不再生成 \(M\)。若某个极小生成组的剩余类线性相关，可从中删去至少一项而保持对剩余向量空间的生成性；刚证明的反向蕴含说明原模仍被其余元素生成，与极小性矛盾。
\end{proof}
\symidx{\(\mu_R(M)\)：局部环上有限生成模的最少生成元个数}

任意生成组的剩余类都必须生成 \(M/\mathfrak mM\)，所以其元素个数至少是 \(\mu_R(M)\)。而提升一组剩余向量空间的基就能达到这个数目。因此在局部环上，不可删去元素的极小生成组也具有最小可能的元素个数；非局部环则可能有不可删去、但个数并非最少的生成组。

\ProseParagraphBreak
例如，设 \(k\) 为域。在 \(R=k[x,y]_{(x,y)}\) 中，\(\mathfrak m=(x,y)R\) 满足
\(\mathfrak m/\mathfrak m^2\cong k^2\)，基为 \(x,y\) 的类。确实一次项给出这两个独立坐标，极大理想平方恰好没有一次项；局部化的分母常数项非零，只会把一次项乘以一个非零标量。因此 \(\mathfrak m\) 需要两个生成元，不是主理想。同样，
\((x,y)^r/(x,y)^{r+1}\) 的基由全次数恰为 \(r\) 的 \(r+1\) 个单项式给出，故 \(\mu_R(\mathfrak m^r)=r+1\)。

\ProseParagraphBreak
极小生成组不一定是自由基。例如 \(R=k[\epsilon]/(\epsilon^2)\) 上的模
\(M=R/(\epsilon)\) 只需一个生成元，但该生成元被非零的 \(\epsilon\) 零化。剩余向量空间的维数测量生成元数，并没有排除原模的关系。上述极小生成组的结论也可对照\cite[Theorem 2.3, pp. 8--9]{Matsumura1989CommutativeRingTheory}。

% !TeX root = ../../Book3.tex
% 纲要标题：### 3.3 Artin 环
% 纲要位置：工作记录/计划纲要.md，第 1991 行。
% 纲要细目（写作范围）：
% * 降链条件与有限长度；通过简单商模证明根幂零及商层的域坐标分解
% * 交换 Artin 环的局部直积分解；满足根幂为零的指数及由极大理想恢复因子
% * 幂零根及零维 Noether 环

\section{Artin 环}
\label{b3:sec:0303}

降链条件对一般模弱于有限长度，但对交换环本身有更强的后果：它迫使极大理想只有有限多个，迫使 Jacobson 根幂零，最终使 \(R\) 作为自身的模具有有限长度。这些性质又把环分成有限个 Artin 局部环。

\subsection{降链条件与有限长度}
\label{b3:subsec:0303:01}

设 \(R\) 为交换含幺环。若非零 \(R\) 模 \(M\) 只有 \(0\) 和 \(M\) 两个子模，则称 \(M\) 为\term[jiandan-mo]{简单模}[simple module]。若 \(R\) 模 \(M\) 存在有限子模链
\[
0=M_0\subsetneq M_1\subsetneq\cdots\subsetneq M_r=M,
\]
且每个 \(M_i/M_{i-1}\) 都简单，则称 \(M\) 为\term[youxian-changdu-mo]{有限长度模}[module of finite length]。上述链称为\term[hecheng-lie]{合成列}[composition series]，各商称为合成因子。第二卷定理12.2.3，即模的 Jordan–Hölder 定理，已经证明其层数与所取链无关，称为模的\term[mo-changdu]{长度}[length of a module]，记作 \(\ell_R(M)\)。第二卷定理12.2.4进一步证明：有限长度等价于同时满足升链与降链条件。对环 \(R\) 而言，这一判据作用于 \(R\) 作为自身的模；任意 Artin 模则未必有限生成。
零模的合成列没有合成因子，长度为零。
\symidx{\(\ell_R(M)\)：模的长度}

\begin{lemma}\label{b3:lem:artin-primes-radical}
设 \(R\ne0\) 为交换含幺 Artin 环。则每个素理想都极大，极大理想只有有限多个，而且 \(J(R)\) 幂零。
\end{lemma}
\begin{proof}
若 \(\mathfrak p\) 素，则整环 \(R/\mathfrak p\) 仍满足降链条件。要证明这个商是域，只需证明其中每个非零元可逆。对这样的元素 \(a\)，主理想链
\((a)\supseteq(a^2)\supseteq\cdots\) 稳定，所以 \(a^n=a^{n+1}b\) 对某个 \(n,b\) 成立。约去非零的 \(a^n\)，得到 \(ab=1\)。该整环是域，故 \(\mathfrak p\) 极大。

若极大理想有无穷多个，从中选取一列两两不同的极大理想 \(\mathfrak m_1,\mathfrak m_2,\ldots\)，令
\(I_n=\mathfrak m_1\cdots\mathfrak m_n\)。不同极大理想之间不能有包含关系，否则较小者的极大性会迫使二者相等。因此对每个 \(i\le n\)，可以选
\(a_i\in\mathfrak m_i\setminus\mathfrak m_{n+1}\)，则
\(\prod_i a_i\in I_n\setminus\mathfrak m_{n+1}\)，因为 \(\mathfrak m_{n+1}\) 素。另一方面
\(I_{n+1}\subseteq I_n\cap\mathfrak m_{n+1}\)，故
\(I_{n+1}\subsetneq I_n\)，与降链条件矛盾。因此极大理想只有
\(\mathfrak m_1,\ldots,\mathfrak m_r\)。

记 \(J=J(R)\)。降链条件给出 \(J^n=J^{n+1}\)，其中 \(n\ge1\)，但稳定并不直接意味着 \(J^n=0\)。此时还不知道 \(J^n\) 有限生成，不能直接对它使用 Nakayama 引理。改看其零化子
\(K=\operatorname{Ann}_R(J^n)\)：证明 \(K=R\) 就等价于证明 \(J^n=0\)，而幂的稳定性给出
\[
(K:J)=\operatorname{Ann}_R(J^{n+1})=K.
\]
若 \(K\ne R\)，严格包含 \(K\) 的理想至少有 \(R\)，降链条件允许从中选一个极小者 \(K'\)。由于 \(K\) 与 \(K'\) 之间没有中间理想，商模 \(Q=K'/K\) 是非零简单模。取 \(x\in K'\setminus K\)，它的像生成 \(Q\)，也就是 \(K'=K+Rx\)，所以 \(Q\) 循环且有限生成。\cref{b3:thm:nakayama}于是给出 \(JQ\ne Q\)；否则这个非零模就会被迫为零。简单性使 \(JQ\) 只能是零或整个 \(Q\)，因此 \(JQ=0\)。这表示 \(Jx\subseteq K\)，于是 \(x\in(K:J)=K\)，与所取 \(x\notin K\) 矛盾。所以 \(K=R\)，即 \(J^n=0\)。
\end{proof}

这个证明用降链条件找到简单商模 \(K'/K\)，把 Nakayama 引理所需的有限生成性落实在这个循环模上。上面的处理可参阅\cite[Theorem 3.2, p. 16]{Matsumura1989CommutativeRingTheory}。

\begin{theorem}\label{b3:thm:artin-finite-length}
设 \(R\) 为交换含幺 Artin 环。则 \(R\) 作为自身的模具有有限长度，因而是 Noether 环；每个有限生成 \(R\) 模也具有有限长度。
\end{theorem}
\begin{proof}
零环的结论直接成立，以下设 \(R\ne0\)。沿用\cref{b3:lem:artin-primes-radical}中极大理想及其幂零交的记号，不同极大理想两两互素，中国剩余定理给出
\[
R/J\cong\prod_{i=1}^r R/\mathfrak m_i.
\]
考虑有限滤过 \(R\supseteq J\supseteq\cdots\supseteq J^n=0\)。对 \(0\le j<n\)，商层 \(L_j=J^j/J^{j+1}\) 被 \(J\) 零化，因而是 \(R/J\) 模；子模和商模又继承降链条件，所以每个 \(L_j\) 是 Artin 模。这是降链条件在商层上提供的有限性，尚未用到环的 Noether 性。

记 \(K_i=R/\mathfrak m_i\)，在 \(R/J\cong\prod_{i=1}^rK_i\) 中，令 \(e_i=(0,\ldots,1,\ldots,0)\)，其中只有第 \(i\) 个坐标为一。由 \(\sum_i e_i=1\) 及 \(e_ie_h=0\)、\(i\ne h\)，有
\[
L_j=e_1L_j\oplus\cdots\oplus e_rL_j.
\]
每个 \(e_iL_j\) 的标量作用只经过第 \(i\) 个域 \(K_i\)，其 \(R\) 子模正是 \(K_i\) 线性子空间。它作为 \(L_j\) 的子模满足降链条件，故对应向量空间也满足降链条件。这样的向量空间必有限维：若含无限线性无关序列，将其扩充为基，逐次删去其中前几个向量所得的线性生成空间就构成严格降链。有限维向量空间中，将基向量逐个加入，就得到相邻商一维的子空间链；各一维商的 \(R\) 子模只有零和自身，所以是简单模。因此各 \(e_iL_j\)，进而每个 \(L_j\)，都具有有限长度。把各商层的合成列取逆像并连接，得到 \(R\) 的有限合成列。

由第二卷定理12.2.4，有限长度推出 Noether 性。若 \(M\) 由 \(s\) 个元素生成，它是 \(R^s\) 的商；有限直和与商都保留有限长度，因此 \(M\) 也具有有限长度。
\end{proof}

\subsection{交换 Artin 环的局部直积分解}
\label{b3:subsec:0303:02}

中国剩余定理已给出一个具体模型：
\[
\mathbb Z/72\mathbb Z\cong\mathbb Z/8\mathbb Z\times\mathbb Z/9\mathbb Z.
\]
右侧两个因子都是局部环，分别对应原环的极大理想 \((\bar2)\)、\((\bar3)\)。剩余类 \(e_1=9\)、\(e_2=64\) 在这个同构下分别对应 \((1,0)\)、\((0,1)\)，所以二者和为一，乘积为零，且各自平方等于自身。任意 \(\mathbb Z/72\mathbb Z\) 模 \(M\) 自动分成 \(9M\oplus64M\)，分别保留 \(2\) 主部分与 \(3\) 主部分，不需要为模另选生成元。一般交换 Artin 环也能由各极大理想恢复这样的局部因子。

\begin{theorem}\label{b3:thm:artin-decomposition}
设 \(R\ne0\) 为交换含幺 Artin 环，极大理想为
\(\mathfrak m_1,\ldots,\mathfrak m_r\)。对任意满足 \(J(R)^n=0\) 的整数 \(n\ge1\)，自然映射给出
\[
R\cong\prod_{i=1}^r R/\mathfrak m_i^n
\cong\prod_{i=1}^r R_{\mathfrak m_i}.
\]
每个因子都是非零 Artin 局部环。相应地，\(1=e_1+\cdots+e_r\) 分解为两两正交的非零幂等元，每个 \(R\) 模都有
\(M=e_1M\oplus\cdots\oplus e_rM\)。这些局部因子由各极大理想唯一确定，分解仅差排列。
\end{theorem}
\begin{proof}
固定满足 \(J(R)^n=0\) 的整数 \(n\ge1\)。要使商映射给出直积分解，需要各 \(\mathfrak m_i^n\) 两两互素，且它们的交为零。对不同的 \(i,j\)，有
\(\mathfrak m_i+\mathfrak m_j=R\)。写 \(a+b=1\)、\(a\in\mathfrak m_i\)、\(b\in\mathfrak m_j\)，展开 \((a+b)^{2n-1}\) 后每一项至少含 \(a^n\) 或 \(b^n\)，故
\(\mathfrak m_i^n+\mathfrak m_j^n=R\)。于是
\[
\bigcap_i\mathfrak m_i^n
=\prod_i\mathfrak m_i^n
=\left(\prod_i\mathfrak m_i\right)^n=J(R)^n=0.
\]
中国剩余定理给出第一个同构。

这些商因子都是局部环，因为包含 \(\mathfrak m_i^n\) 的极大理想必包含 \(\mathfrak m_i\)：对每个 \(a\in\mathfrak m_i\)，有 \(a^n\in\mathfrak m_i^n\)，再用素性即可。因 \(\mathfrak m_i\) 本身极大，这样的极大理想只能是 \(\mathfrak m_i\)。又因 \(\mathfrak m_i^n\subseteq\mathfrak m_i\subsetneq R\)，商环非零；它继承 Artin 性，所以是非零 Artin 局部环。

在上述直积中，令 \(e_i\) 对应第 \(i\) 个坐标为 \(1\)、其余为 \(0\) 的元素。它们非零、两两正交且和为 \(1\)。对每个 \(R\) 模 \(M\)，等式 \(m=\sum_i e_i m\) 给出所述分解；若 \(\sum_i m_i=0\)、\(m_i\in e_iM\)，逐次乘 \(e_j\) 得 \(m_j=0\)，故和是直和。

局部因子还可直接从原环的极大理想恢复，这将使它们不依赖所取的 \(n\)。考虑任意有限直积分解
\[
R\cong\prod_{j=1}^s B_j,
\]
其中 \(B_j\) 是非零局部环，唯一极大理想记作 \(\mathfrak n_j\)。沿此同构识别 \(R\) 与直积。若 \(\mathfrak m\) 是 \(R\) 的极大理想，各坐标幂等元在域 \(R/\mathfrak m\) 中只能取 \(0\) 或 \(1\)，且由正交性及其和为 \(1\)，恰有一个坐标幂等元取值 \(1\)。设它对应第 \(j\) 个坐标，则其他坐标因子都在 \(\mathfrak m\) 中，而第 \(j\) 个坐标上的商必须是域。因此
\[
\mathfrak m=\mathfrak m_j
=B_1\times\cdots\times\mathfrak n_j\times\cdots\times B_s.
\]
反过来，每个 \(\mathfrak m_j\) 的商是域 \(B_j/\mathfrak n_j\)，故它们恰好穷尽 \(R\) 的极大理想。

第 \(j\) 个坐标投影将 \(R\setminus\mathfrak m_j\) 映到 \(B_j\) 的单位，因而诱导映射
\[
R_{\mathfrak m_j}\longrightarrow B_j,
\qquad a/u\longmapsto a_j u_j^{-1}.
\]
它是满射，因为坐标投影已经满射。若其值为零，则 \(a_j=0\)；第 \(j\) 个坐标幂等元 \(e_j\notin\mathfrak m_j\) 满足 \(e_j a=0\)，所以 \(a/u=0\)。映射的核为零，从而 \(R_{\mathfrak m_j}\cong B_j\)。应用于中国剩余定理所得的分解，就得到 \(R_{\mathfrak m_i}\cong R/\mathfrak m_i^n\)。对任意另一组非零局部因子，极大理想集合仍恢复同一组局部化，故因子只差排列和同构。
\end{proof}

零环没有极大理想，采用空直积为零环的约定后，它对应空的局部分解。零环上的幺性模只有零模，相应的模分解也是空直和。

\subsection{幂零根与零维 Noether 环}
\label{b3:subsec:0303:03}

本小节把“零维”用于每个素理想都极大的非零环；一般 Krull 维数的定义和基本定理留在\chapref{b3:ch:09}。零环作为单独的平凡情形处理。

\begin{theorem}\label{b3:thm:noether-dimension-zero}
设 \(R\) 为非零交换含幺环。下列条件等价：
\begin{equivalences}
\item \(R\) 为 Artin 环；
\item \(R\) 为 Noether 环，且每个素理想都极大。
\end{equivalences}
此时幂零根与 Jacobson 根相等，均为幂零理想。\(R\) 约化当且仅当它是有限个域的直积。
\end{theorem}
\begin{proof}
由\cref{b3:lem:artin-primes-radical}及\cref{b3:thm:artin-finite-length}，第一项推出第二项。

反过来，设 \(R\) 为 Noether 环且每个素理想极大。第一章的\cref{b3:prop:noether-spectrum-components}表明其极小素理想只有有限多个，且每个素理想包含一个极小素理想。现在包含在两个素理想之间只能是相等，所以整个素谱就是有限个极大理想
\(\mathfrak m_1,\ldots,\mathfrak m_r\)。由\cref{b3:thm:radical-prime-intersection}，
\[
\sqrt{(0)}=\bigcap_i\mathfrak m_i=J(R).
\]
Noether 条件使幂零根有限生成；它的每个元素幂零，故由\cref{b3:prop:finite-nil-ideal}，该理想的某个幂为零。

记 \(J=J(R)\)，中国剩余定理仍给出 \(R/J\cong\prod_i R/\mathfrak m_i\)。这里各理想 \(J^j\) 因 Noether 性而有限生成，所以商层 \(J^j/J^{j+1}\) 是有限生成 \(R\) 模。它被 \(J\) 零化，因而也是有限生成的 \(R/J\) 模。与前面从降链条件得到 Artin 商层不同，这个方向依靠有限生成性：按坐标幂等元分解以后，每个域上的向量空间分量由原来有限个生成元的像生成，故有限维。因此各商层具有有限长度；沿有限的 \(J\) 幂滤过连接合成列，得到 \(R\) 具有有限长度，故满足降链条件。

最后，若 \(R\) 约化，上述公共根为零，直接得到 \(R\cong\prod_i R/\mathfrak m_i\)。有限个域的直积没有非零幂零元，反向也成立。
\end{proof}

上述两种有限性给出不同的蕴含关系。在 \(R\) 为 Artin 环的假设下，有
\[
\left.
\begin{gathered}
\#\operatorname{Max}R<\infty,\quad J^n=0,\\
J^i/J^{i+1}\;\text{满足降链条件}
\end{gathered}
\right\}
\Longrightarrow\ell_R(R)<\infty
\Longrightarrow R\;\text{为 Noether 环}.
\]
而在 \(R\) 为 Noether 环且每个素理想都极大的假设下，有
\[
\left.
\begin{gathered}
\operatorname{Spec}R\;\text{有限},\quad J=\sqrt{(0)},\quad J^n=0,\\
J^i/J^{i+1}\;\text{有限生成}
\end{gathered}
\right\}
\Longrightarrow\ell_R(R)<\infty
\Longrightarrow R\;\text{为 Artin 环}.
\]
\ProseParagraphBreak
去掉 Noether 条件，零维便不足以推出 Artin 性。\secref{b3:sec:0102}中的无限域直积 \(R=\prod_{n\ge1}k\) 已经给出反例：它的每个素理想都极大，而使前 \(n\) 个坐标为零的理想构成严格降链。

% !TeX root = ../../Book3.tex
% 纲要标题：### 3.4 有限性在变换下的行为
% 纲要位置：工作记录/计划纲要.md，第 1897 行。
% 纲要细目（写作范围）：
% * 商、局部化和有限代数
% * Noether 性并非任意子环继承
% * 有限性假设的常见误用
% ---

\section{有限性在变换下的行为}
\label{b3:sec:0304}

商、局部化和有限代数都保留相当多的有限性，而取子环一般不保留。把这些结论放在一起，可以避免在具体计算中因“环变小了”或“只剩一个点了”就误判其理想结构。本节还区分茎上的条件与基本开邻域上的条件。

\subsection{商、局部化与有限代数}
\label{b3:subsec:0304:01}

\begin{proposition}\label{b3:prop:chain-conditions-localization}
设 \(R\) 为交换含幺环，\(I\subseteq R\) 为理想，\(S\subseteq R\) 为含一的乘法集。若 \(R\) 为 Noether 环，则 \(R/I\) 与 \(S^{-1}R\) 仍为 Noether 环；若 \(R\) 为 Artin 环，则二者仍为 Artin 环。对任意 \(R\) 模 \(M\)，若 \(M\) 有限生成或有限展示，则 \(S^{-1}M\) 作为 \(S^{-1}R\) 模仍分别满足相应条件。
\end{proposition}
\begin{proof}
商环的理想链取逆像，严格包含仍严格，故两种链条件都传给商。局部化环中每个理想都是其收缩的扩张，见\cref{b3:prop:ideal-localization-saturation}。因此一条严格理想链收缩后仍严格；原环的相应链条件排除了无限严格升链或降链。

模的生成元取分式 \(m_i/1\) 就生成局部化模。对有限展示
\(R^a\rightarrow R^b\rightarrow M\rightarrow0\)，使用\cref{b3:thm:localization-exact}，得到
\((S^{-1}R)^a\rightarrow(S^{-1}R)^b\rightarrow S^{-1}M\rightarrow0\)，故有限展示也被保持。
\end{proof}

\begin{proposition}\label{b3:prop:finite-algebra-noether}
设 \(R\) 为交换含幺环，\(A\) 为交换含幺 \(R\) 代数。若 \(A\) 作为 \(R\) 模有限生成，则称 \(A\) 为有限 \(R\) 代数。设 \(A\) 满足这一条件。
\begin{enumerate}
\item 若 \(R\) 为 Noether 环，则 \(A\) 为 Noether 环。
\item 若 \(R\) 为 Artin 环，则 \(A\) 为 Artin 环。
\item 若 \(M\) 为有限生成 \(A\) 模，则它也是有限生成 \(R\) 模。
\end{enumerate}
\end{proposition}
\begin{proof}
第一项中，\(A\) 作为有限生成 \(R\) 模满足升链条件，而每个 \(A\) 理想都是 \(R\) 子模，所以理想升链稳定。第二项同理：由\cref{b3:thm:artin-finite-length}，\(A\) 作为 \(R\) 模具有有限长度，其 \(R\) 子模降链稳定，特别地 \(A\) 理想降链稳定。第三项中，若 \(a_1,\ldots,a_s\) 生成 \(A\) 作为 \(R\) 模，\(m_1,\ldots,m_t\) 生成 \(M\) 作为 \(A\) 模，则全部 \(a_i m_j\) 生成 \(M\) 作为 \(R\) 模。
\end{proof}

这里的“有限”比有限型严格得多。\(k[x]\) 是域 \(k\) 上的有限型代数，却不是 Artin 环，因为 \((x)\supsetneq(x^2)\supsetneq\cdots\)。Noether 底环上的有限型代数保持 Noether 性，依据是 Hilbert 基定理；Artin 底环上的有限代数保持 Artin 性，依据则是模的有限长度。

\begin{proposition}\label{b3:prop:finite-presentation-open-cover}
设 \(R\) 为交换含幺环，\(M\) 为 \(R\) 模，\(f_1,\ldots,f_r\in R\)（\(r\ge1\)）生成单位理想。若每个 \(M_{f_i}\) 是有限展示
\(R_{f_i}\) 模，则 \(M\) 有限展示。若每个 \(R_{f_i}\) 是 Noether 环，则 \(R\) 为 Noether 环。
\end{proposition}
\begin{proof}
由\cref{b3:prop:principal-cover-finiteness}，\(M\) 有限生成，取满射
\(R^n\rightarrow M\)，核为 \(K\)。局部化后仍正合，且目标有限展示，所以由
\cref{b3:prop:finite-presentation-kernel}，每个 \(K_{f_i}\) 有限生成。再次使用有限基本开覆盖上的有限生成性判据，得到 \(K\) 有限生成，因而 \(M\) 有限展示。

对第二句，任取理想 \(I\subseteq R\)。其局部化是 Noether 环 \(R_{f_i}\) 的理想，故有限生成。对 \(I\) 使用同一有限生成性判据，得到每个理想都有限生成，所以 \(R\) 为 Noether 环。
\end{proof}

\subsection{Noether 环的非 Noether 子环}
\label{b3:subsec:0304:02}

设 \(k\) 为域，取多项式环 \(B=k[x,y]\) 的子环
\[
A=k+xk[x,y]
=k[x,xy,xy^2,\ldots].
\]
等号右侧表示常数项之外的每个单项式都至少含一个 \(x\)；它对乘法封闭，并且与 \(B\) 具有相同的单位元。环 \(B\) 由 Hilbert 基定理是 Noether 环，\(A\) 却不是。

\ProseParagraphBreak
为此令 \(I=xk[x,y]\)，它是 \(A\) 的理想，\(A/I\cong k\)，而
\(I^2=x^2k[x,y]\)。于是
\[
I/I^2
=\bigoplus_{j\ge0}k\cdot(xy^j+I^2)
\]
为无限维 \(k\) 向量空间。如果 \(I\) 由有限个元素作为 \(A\) 理想生成，模掉 \(I^2\) 后，这些元素的剩余类就会在 \(A/I=k\) 上生成 \(I/I^2\)，与无限维性矛盾。故 \(I\) 不有限生成，\(A\) 不是 Noether 环。

\ProseParagraphBreak
这个例子也显示问题所在：在较大的环 \(B\) 中，\(I\) 由一个元素 \(x\) 生成，因为可以任意乘以 \(y\)；在 \(A\) 中没有这个乘子。取子环虽然减少了元素，却也减少了生成理想时允许使用的系数，可能使有限生成性失败。

\subsection{有限性假设的适用范围}
\label{b3:subsec:0304:03}

有限个生成元只描述模本身，并不保证它的所有子模也能如此描述。设 \(k\) 为域，环 \(k[x_1,x_2,\ldots]\) 作为自身的模由 \(1\) 生成，其理想 \((x_1,x_2,\ldots)\) 却不有限生成：任意有限组候选生成元只涉及有限多个变量，令这些变量全为零后，它们生成的理想也映为零，而其余变量的像非零。在 Noether 底环上，\cref{b3:prop:noether-finite-modules}保证有限生成模的每个子模都有限生成。

\ProseParagraphBreak
即使每个点处都能找到有限组生成元，也未必已经得到有限的全局描述。\subsecref{b3:subsec:0204:03}中的
\(\bigoplus_p\mathbb Z/p\mathbb Z\) 在每个素理想处都有限生成，但全局不是有限生成模。\cref{b3:prop:principal-cover-finiteness}及\cref{b3:prop:finite-presentation-open-cover}要求的是一组基本开邻域上的有限数据；仅检查各个茎不足以自动得到这些邻域。

\ProseParagraphBreak
环的 Noether 性也有相同的区别。\secref{b3:sec:0102}中的无限域直积 \(R=\prod_{n\ge1}k\) 满足：每个元素 \(a\) 都可写成 \(a=a^2b\)。若 \(a\in\mathfrak p\)，则 \(1-ab\notin\mathfrak p\)，而 \(a(1-ab)=0\)，所以 \(a/1=0\) 于 \(R_{\mathfrak p}\)。于是 \(\mathfrak pR_{\mathfrak p}=0\)，每个局部环 \(R_{\mathfrak p}\) 都是域。然而，令 \(e_n\) 为第 \(n\) 个坐标为一、其余坐标为零的元素，理想 \(Re_1+\cdots+Re_n\) 组成严格升链，故 \(R\) 不是 Noether 环。各点处的 Noether 性也不能代替基本开邻域上的条件。

\ProseParagraphBreak
降链条件也不能一般地保证模有限生成。设 \(p\) 为素数，考虑整数模
\[
C_{p^\infty}=\mathbb Z[1/p]/\mathbb Z
\]
记 \(C_n=\langle p^{-n}+\mathbb Z\rangle\)（\(n\ge0\)），则 \(C_0=0\)，\(C_n\) 是阶为 \(p^n\) 的循环群，且
\[
C_0\subsetneq C_1\subsetneq C_2\subsetneq\cdots,
\qquad C_{p^\infty}=\bigcup_{n\ge0}C_n.
\]
任意阶为 \(p^n\) 的元素都能写成 \(a/p^n+\mathbb Z\)，其中 \(p\nmid a\)。取整数 \(u,v\) 使 \(ua+vp^n=1\)，就有
\(u(a/p^n+\mathbb Z)=p^{-n}+\mathbb Z\)，所以该元素恰好生成 \(C_n\)。此外，被 \(p^n\) 零化的元素恰为 \(C_n\) 的元素。故任意阶为 \(p^n\) 的子群都包含于 \(C_n\)，比较元素个数可知它等于 \(C_n\)：每一阶的子群只有一个。

\ProseParagraphBreak
现在设 \(H\) 为子模。若 \(H\) 中元素的阶无界，则对每个 \(n\) 都能取到阶为 \(p^r\)、\(r\ge n\) 的元素；它生成 \(C_r\)，因而 \(C_n\subseteq H\)，于是 \(H=C_{p^\infty}\)。若阶有界，则 \(H\subseteq C_n\) 对某个 \(n\) 成立。在有限群 \(H\) 中选阶最大的元素，设其阶为 \(p^r\)，则它生成 \(C_r\)，而 \(H\) 中其余元素都被 \(p^r\) 零化，所以 \(H=C_r\)。这就证明真子模恰为各个有限循环群 \(C_r\)。任何子模降链一旦离开全模便进入有限群，此后稳定，故 \(C_{p^\infty}\) 是 Artin 模；任意有限组元素都落在某个 \(C_n\) 中，不能生成全模，故它不是有限生成模。

\ProseParagraphBreak
\cref{b3:tab:finiteness-counterexamples}汇集了这些反例，以及保证相应结论成立的充分条件。

\begin{table}[htbp]
\centering
\caption{有限性推论的反例与适用条件}
\label{b3:tab:finiteness-counterexamples}
\begin{tabularx}{\textwidth}{@{}>{\raggedright\arraybackslash}p{.25\textwidth}>{\raggedright\arraybackslash}p{.30\textwidth}>{\raggedright\arraybackslash}X@{}}
\toprule
一般不成立的推论 & 反例 & 保证相应结论的条件\\
\midrule
有限生成模的子模都有限生成 & 无限变量多项式环的变量理想 & 底环是 Noether 环\\
各个素理想处有限生成，便全局有限生成 & \(\bigoplus_{p\text{ 素数}}\mathbb Z/p\mathbb Z\) & 一组覆盖素谱的基本开邻域上有限生成\\
Artin 模必有限生成 & 整数模 \(C_{p^\infty}\) & 模同时满足 Noether 条件，此时具有有限长度\\
\bottomrule
\end{tabularx}
\end{table}

有限生成使有限组生成元可以共用分母，也是 Nakayama 引理对模的要求；有限展示进一步要求关系模有限生成。Noether 条件控制所有子模的有限生成性，Artin 条件控制降链；对交换环本身，后者还给出\cref{b3:thm:artin-decomposition}中的局部直积分解。


\begin{chaptersummary}
有限展示要求模有有限个生成元，且全部关系由有限多个基本关系线性生成，并不把关系总数限制为有限。Noether 环使每个有限生成模的子模仍有限生成。Hilbert 基定理保证 Noether 环上的有限型代数仍是 Noether 环，因而其中的理想与有限矩阵的关系模具有有限描述。

\ProseParagraphBreak
行列式技巧不要求生成元线性无关，却能产生一个统一的首一零化方程。对恒等映射使用它，再利用 Jacobson 根的单位判据，便得到 Nakayama 引理。在局部环上，极小生成元数因此等于剩余向量空间的维数；这项维数不负责判断关系是否消失。

\ProseParagraphBreak
交换 Artin 环具有有限长度，并分解为各极大理想处的局部环。反向刻画需要 Noether 性与零维同时成立。有限代数、商及局部化保留适当的链条件，任意子环却不保留；后续整扩张与维数论中的有限性假设将继续依赖这些区别。
\end{chaptersummary}

\BookExercises

\begin{exercise}\label{b3:ex:03-cyclic-presentation}
设 \(k\) 为域，\(R=k[x,y]\)、\(\mathfrak m=(x,y)\)、\(I=(x^2,xy,y^2)\)。给出循环模 \(R/I\) 的有限展示，并说明三个基本关系 \(x^2\cdot(1+I)=xy\cdot(1+I)=y^2\cdot(1+I)=0\) 都不能直接删去。再通过 \(I/\mathfrak mI\) 证明：针对生成元 \(1+I\)，其关系模不能由两个元素生成。
\end{exercise}
\begin{solution}
关系模是商映射 \(R\to R/I\) 的核 \(I\)，故有有限展示
\[
R^3\xrightarrow{\begin{pmatrix}x^2&xy&y^2\end{pmatrix}}R
\longrightarrow R/I\longrightarrow0.
\]
一个模生成元为 \(1+I\)，全部关系是这三个基本关系的 \(R\) 线性组合。若删去某个二次单项式，它不能由另两个生成：乘以多项式后的二次项只能是这两个单项式的 \(k\) 线性组合。

更一般地，\(\mathfrak mI=(x,y)^3\) 只含全次数至少为三的项，所以 \(I/\mathfrak mI\) 以 \(x^2,xy,y^2\) 的类为 \(k\) 基，维数为三。由于 \(\mathfrak m\) 在这个商上作用为零，\(R\) 的作用经过 \(R/\mathfrak m=k\)。若任意两个元素生成 \(I\)，它们的剩余类就生成该三维向量空间，矛盾。因此不是换两条巧选的基本关系就能代替这三条。这里计算的是关系模的生成元数，模 \(R/I\) 本身仍只有一个生成元。
\end{solution}

\begin{exercise}\label{b3:ex:03-two-variable-relation}
设 \(k\) 为域，令 \(R=k[x,y]\)。求映射 \(R^2\rightarrow(x,y)\)、\((a,b)\mapsto ax+by\) 的核，并给出 \((x,y)\) 的有限展示。
\end{exercise}
\begin{solution}
若 \(ax+by=0\)，模掉 \(y\) 后在 \(k[x]\) 中有 \(\bar a x=0\)，故 \(a=yc\)。代回得到 \(y(cx+b)=0\)，整环性给出 \(b=-xc\)。因此核由 \((y,-x)\) 生成，得到正合列
\[
0\longrightarrow R\xrightarrow{c\mapsto(yc,-xc)}R^2
\longrightarrow(x,y)\longrightarrow0.
\]
第一个映射单射，因为 \(y\ne0\) 且 \(R\) 为整环。
\end{solution}

\begin{exercise}\label{b3:ex:03-finite-extension-modules}
设 \(R\) 为交换含幺环，\(0\rightarrow L\rightarrow M\rightarrow N\rightarrow0\) 为 \(R\) 模的正合列。证明若 \(L,N\) 有限生成，则 \(M\) 有限生成；同时给出 \(M\) 有限生成而 \(L\) 不有限生成的例子。
\end{exercise}
\begin{solution}
取 \(L\) 的有限生成元，并提升 \(N\) 的有限生成元到 \(M\)。对任意 \(m\in M\)，先减去提升元的适当线性组合，使差落入 \(L\)，再用 \(L\) 的生成元表示，故这些有限个元素生成 \(M\)。
反向取域 \(k\)，令 \(R=k[x_1,x_2,\ldots]\)、\(M=R\)、\(L=(x_1,x_2,\ldots)\)。\(M\) 循环，但 \(L\) 不有限生成：任意有限组候选生成元只涉及有限个变量，将这些变量置零后，它们全变为零，而一个未出现的变量仍非零。此处底环不是 Noether 环。
\end{solution}

\begin{exercise}\label{b3:ex:03-cusp-finite}
设 \(k\) 为域，令 \(A=k[t^2,t^3]\subseteq k[t]\)、\(R=k[t^2]\)。证明 \(A\) 是有限 \(R\) 代数并为 Noether 环；比较 \(k[t,y]\) 作为 \(k[t]\) 代数的有限型性与有限性。
\end{exercise}
\begin{solution}
每个 \(A\) 中的单项式次数都为偶数，或为至少 \(3\) 的奇数，故
\(A=R+Rt^3\)。所以 \(A\) 作为 \(R\) 模由 \(1,t^3\) 生成。\(R\) 同构于一元多项式环，是 Noether 环，故由\cref{b3:prop:finite-algebra-noether}，\(A\) 为 Noether 环。
另一方面 \(k[t,y]\) 作为 \(k[t]\) 代数由 \(y\) 生成，是有限型代数；若它作为模由有限个多项式生成，生成元的 \(y\) 次数有共同上界，而 \(k[t]\) 系数的组合不能超过这个上界，不能得到任意高次的 \(y^n\)。所以它不是有限代数。
\end{solution}

\begin{exercise}\label{b3:ex:03-determinant-integrality}
设 \(R\) 为整环，\(K\) 为其分式域，\(0\ne M\subseteq K\) 为有限生成 \(R\) 子模。若 \(z\in K\) 满足 \(zM\subseteq M\)，证明 \(z\) 满足一个 \(R\) 系数的首一多项式。
\end{exercise}
\begin{solution}
乘 \(z\) 给出 \(M\) 的 \(R\) 线性自同态 \(m\mapsto zm\)，其中“自同态”指从 \(M\) 到自身的同态，不要求可逆。对 \(I=R\) 应用\cref{b3:thm:determinant-trick}，得到首一 \(P(T)\in R[T]\)，使
\(P(z)\cdot m=0\) 对所有 \(m\in M\) 成立。取一个非零 \(m\in M\subseteq K\)，在域中约去它，得到 \(P(z)=0\)。非零子模的条件负责最后的约去；若 \(M=0\)，这个论证不能给 \(z\) 任何限制。
\end{solution}

\begin{exercise}\label{b3:ex:03-general-nakayama-factor}
在 \(R=\mathbb Z\)、\(I=(2)\)、\(M=\mathbb Z/3\mathbb Z\) 中，找出 \(a\in I\) 使 \((1-a)M=0\)。解释为什么这不意味着 \(M=0\)。
\end{exercise}
\begin{solution}
取 \(a=4\)，则 \(1-a=-3\) 零化 \(M\)，而 \(4\in(2)\)。这里 \(-3\) 不是整数环的单位，不能从 \((-3)M=0\) 约去 \(-3\)。\((2)\) 也不包含于
\(J(\mathbb Z)=0\)，因为所有极大理想 \((p)\) 的交为零。所以满足的是行列式技巧给出的统一零化元结论，没有满足 Nakayama 的 Jacobson 根假设。
\end{solution}

\begin{exercise}\label{b3:ex:03-local-matrix}
设 \(R\) 为交换含幺局部环，极大理想为 \(\mathfrak m\)，剩余域为 \(\kappa=R/\mathfrak m\)，\(n\ge1\)，\(A\in M_n(R)\)。证明 \(A\) 可逆当且仅当其模 \(\mathfrak m\) 的矩阵 \(\bar A\in M_n(\kappa)\) 可逆。再设 \(M\) 为有限生成 \(R\) 模，\(m_1,\ldots,m_n\) 与 \(m'_1,\ldots,m'_n\) 是它的两组极小生成元。证明任意满足
\[
m'_i=\sum_{j=1}^n b_{ij}m_j\qquad(1\le i\le n)
\]
的转换矩阵 \(B=(b_{ij})\in M_n(R)\) 都可逆。
\end{exercise}
\begin{solution}
若 \(A\) 可逆，取商后仍可逆。反向若 \(\bar A\) 可逆，则
\(\det A\bmod\mathfrak m=\det\bar A\ne0\)，故 \(\det A\) 为 \(R\) 的单位，伴随矩阵公式给出 \(A^{-1}=(\det A)^{-1}\operatorname{adj}(A)\)。
由\cref{b3:cor:local-generators}，两组极小生成元模掉 \(\mathfrak mM\) 后都是 \(M/\mathfrak mM\) 的 \(\kappa\) 基。因此上述等式约化后给出两组基之间的转换矩阵 \(\bar B\)，它必可逆；由第一部分，\(B\) 在 \(R\) 上可逆。生成元未必是自由基，所以表示它们的矩阵未必唯一；这里证明的是每个这样的矩阵都可逆。
\end{solution}

\begin{exercise}\label{b3:ex:03-minimal-power-generators}
设 \(k\) 为域，\(d\ge1\)，令 \(R=k[x_1,\ldots,x_d]_{(x_1,\ldots,x_d)}\)、\(\mathfrak m=(x_1,\ldots,x_d)R\)。求 \(\mu_R(\mathfrak m^r)\)、\(r\ge0\)。
\end{exercise}
\begin{solution}
\(\mathfrak m^r/\mathfrak m^{r+1}\) 以全次数 \(r\) 的单项式为 \(k\) 基。局部化分母的常数项非零，在这个商中只影响其非零常数倍，不引入新的线性关系。因此
\[
\mu_R(\mathfrak m^r)=\dim_k(\mathfrak m^r/\mathfrak m^{r+1})
=\binom{r+d-1}{d-1}.
\]
末个数是在 \(d\) 个非负整数中分配总和 \(r\) 的方案数。\(r=0\) 时即 \(R/\mathfrak m\) 的维数 \(1\)，与 \(R\) 作为自身的模循环一致。
\end{solution}

\begin{exercise}\label{b3:ex:03-nonlocal-minimal}
证明整数模 \(\mathbb Z\) 同时有极小生成组 \(\{1\}\) 和 \(\{2,3\}\)。说明这不违反局部环上的极小生成元结论。
\end{exercise}
\begin{solution}
\(\{1\}\) 显然生成；\(\{2,3\}\) 也生成，因为 \(3-2=1\)，但删去任意一个只得到 \(2\mathbb Z\) 或 \(3\mathbb Z\)，所以它也极小。“极小”是不能删去元素，并不是在所有生成组中个数最少。
整数环不是局部环，因此不能把不同极大理想的剩余信息统一成同一个剩余向量空间。局部结论所用的 \(\mathfrak m\) 唯一性在此缺失。
\end{solution}

\begin{exercise}\label{b3:ex:03-surjective-endomorphism}
设 \(R\) 为交换含幺环，\(M,N\) 为 \(R\) 模，其中 \(M\) 有限生成。若 \(M\cong M\oplus N\) 为 \(R\) 模同构，证明 \(N=0\)。再用整数模 \(\mathbb Z\) 的一个单自同态说明“满自同态可逆”不能改成“单自同态可逆”。
\end{exercise}
\begin{solution}
将给定同构 \(M\rightarrow M\oplus N\) 与第一投影复合，得到 \(M\) 的满自同态。由\cref{b3:cor:finite-surjective-endomorphism}它单射，而其核在原同构下恰为 \(0\oplus N\)，故 \(N=0\)。
整数模的乘二映射为单射但不是满射，因 \(1\) 没有原像，所以不具有逆映射。
\end{solution}

\begin{exercise}\label{b3:ex:03-idempotent-ideal}
设 \(R\) 为交换含幺环，理想 \(I\subseteq R\) 有限生成且 \(I^2=I\)。证明存在幂等元 \(e\in R\) 使 \(I=eR\)。若 \(R\) 局部，说明只能有 \(I=0\) 或 \(I=R\)。
\end{exercise}
\begin{solution}
把 \(I\) 看作有限生成模，行列式技巧应用于 \(I=II\)，得到 \(e\in I\) 使
\((1-e)I=0\)。于是 \(x=ex\) 对所有 \(x\in I\) 成立，故 \(I\subseteq eR\)；另一包含由 \(e\in I\) 得到。再取 \(x=e\)，便有 \(e=e^2\)。
在局部环中，\(e\) 与 \(1-e\) 至少一个是单位，因为二者之和为 \(1\)。由
\(e(1-e)=0\)，若 \(e\) 可逆则 \(e=1\)，若 \(1-e\) 可逆则 \(e=0\)。因此 \(I\) 只能是所述两个理想。
\end{solution}

\begin{exercise}\label{b3:ex:03-integer-artin-product}
显式分解 \(\mathbb Z/360\mathbb Z\) 为 Artin 局部环的直积，并找出三个对应的正交幂等元。
\end{exercise}
\begin{solution}
\(360=8\cdot9\cdot5\)，三因子两两互素，所以
\[
\mathbb Z/360\mathbb Z\cong
\mathbb Z/8\mathbb Z\times\mathbb Z/9\mathbb Z\times\mathbb Z/5\mathbb Z.
\]
可取 \(e_1=225\)、\(e_2=280\)、\(e_3=216\) 的剩余类。它们在对应因子中取 \(1\)，在另外两个因子中取 \(0\)；由中国剩余同构可知它们幂等、两两正交，且和为 \(1\)。也可直接核对 \(225+280+216=721\equiv1\pmod{360}\)。
\end{solution}

\begin{exercise}\label{b3:ex:03-artin-square-example}
设 \(k\) 为域，对 \(A=k[x,y]/(x^2,y^2)\)，求唯一极大理想、其幂零指数、环的长度，以及零化该极大理想的元素组成的理想。
\end{exercise}
\begin{solution}
\(A\) 的 \(k\) 基为 \(1,\bar x,\bar y,\bar x\bar y\)。非零常数项的元素可用有限几何级数求逆，零常数项的元素幂零，所以唯一极大理想为
\(\mathfrak m=(\bar x,\bar y)\)。有
\(\mathfrak m^2=k\bar x\bar y\ne0\)、\(\mathfrak m^3=0\)，幂零指数为 \(3\)。
商层 \(A/\mathfrak m\)、\(\mathfrak m/\mathfrak m^2\)、\(\mathfrak m^2\) 的维数分别为 \(1,2,1\)，将中间二维层细分即得长度 \(4\)。
若 \(a+b\bar x+c\bar y+d\bar x\bar y\) 被 \(\bar x,\bar y\) 同时零化，分别相乘得到 \(a=c=0\) 与 \(a=b=0\)，故零化子为 \(k\bar x\bar y\)。
\end{solution}

\begin{exercise}\label{b3:ex:03-polynomial-artin-product}
设 \(k\) 为域。分解 \(A=k[t]/\bigl(t^2(t-1)^3\bigr)\) 为局部因子，求其长度及约化商。
\end{exercise}
\begin{solution}
\((t^2)\) 与 \(\bigl((t-1)^3\bigr)\) 互素，所以
\[
A\cong k[t]/(t^2)\times k[t]/\bigl((t-1)^3\bigr).
\]
两个因子分别是以 \((\bar t)\) 与 \((\overline{t-1})\) 为唯一极大理想的 Artin 局部环，具有长度 \(2,3\)，因此 \(A\) 的长度为 \(5\)。幂零根在两个因子中分别为这两个极大理想，约化商为 \(k\times k\)。
\end{solution}

\begin{exercise}\label{b3:ex:03-nilpotent-maximal-not-artin}
设 \(k\) 为域，\(V\) 为无限维 \(k\) 向量空间，在 \(R=k\oplus V\) 上规定
\((a,u)(b,v)=(ab,av+bu)\)。证明它是局部环，极大理想平方为零，但它不是 Artin 环。
\end{exercise}
\begin{solution}
该乘法直接满足结合律及分配律，单位为 \((1,0)\)。若 \(a\ne0\)，
\((a,u)^{-1}=(a^{-1},-a^{-2}u)\)；若 \(a=0\)，元素平方为零，所以非单位集恰为 \(\mathfrak m=0\oplus V\)，且 \(\mathfrak m^2=0\)。
\(V\) 的每个 \(k\) 子空间都是 \(R\) 理想，因为 \(V\) 对它作用为零，标量 \(k\) 按通常方式作用。选一列线性无关向量并扩充为基，逐次删去前面的向量所得子空间构成严格降链。因此 \(R\) 不是 Artin 环。极大理想幂零本身不能替代有限性。
\end{solution}

\begin{exercise}\label{b3:ex:03-zero-dimensional-reduced}
设 \(k\) 为域。证明有限维约化交换含幺 \(k\) 代数是有限个有限域扩张的直积（零代数对应空乘积）。指出当 \(k\) 代数闭时这些因子是什么。
\end{exercise}
\begin{solution}
设该代数为 \(A\)。若 \(A=0\)，按约定它就是空乘积。以下设 \(A\ne0\)。有限维性使理想升链、降链都稳定，故 \(A\) 为 Artin 环；约化性结合
\cref{b3:thm:noether-dimension-zero}给出有限个域的直积。各域为原代数的商，因此作为 \(k\) 向量空间仍有限维，是有限域扩张。若 \(k\) 代数闭，每个有限扩张中的元素都有 \(k\) 上的最小多项式，该不可约多项式只能为一次，故各因子均等于 \(k\)。
\end{solution}

\begin{exercise}\label{b3:ex:03-artin-local-length}
设 \(R\) 为交换含幺 Artin 局部环，极大理想为 \(\mathfrak m\)，剩余域为 \(\kappa=R/\mathfrak m\)，\(M\) 为有限生成 \(R\) 模。证明
\[
\ell_R(M)=\sum_{i\ge0}
\dim_\kappa(\mathfrak m^iM/\mathfrak m^{i+1}M).
\]
推导 \(\ell_R(M)\ge\mu_R(M)\)，并刻画等号成立的情形。
\end{exercise}
\begin{solution}
由\cref{b3:thm:artin-finite-length}，\(R\) 是 Noether 环；再由\cref{b3:prop:noether-finite-modules}，有限生成模 \(M\) 的每个子模 \(\mathfrak m^iM\) 都有限生成。因此各商层也是有限生成 \(R\) 模；它们被 \(\mathfrak m\) 零化，遂成为有限维 \(\kappa\) 向量空间。又因 \(\mathfrak m\) 幂零，和中只有有限个非零项。按各层一组基的前若干项构造向量空间旗标。这些商层上的 \(R\) 作用都经过 \(R/\mathfrak m=\kappa\)，所以其 \(R\) 子模恰是 \(\kappa\) 线性子空间；每个一维商只有零子模与自身，因而是简单 \(R\) 模 \(\kappa\)。将各层的旗标取逆像并连接，就得到总层数为右侧的合成列。
第一项为 \(\dim_\kappa M/\mathfrak mM=\mu_R(M)\)，其他项非负，因此有不等式。若等号成立，则 \(\mathfrak mM/\mathfrak m^2M=0\)。上面已经证明 \(\mathfrak mM\) 有限生成，对它应用\cref{b3:thm:nakayama}，得到 \(\mathfrak mM=0\)。反向若 \(\mathfrak mM=0\)，只有第一项非零，当然取等号。
\end{solution}

\begin{exercise}\label{b3:ex:03-artin-localized-length}
设 \(R\) 为交换含幺 Artin 环，极大理想为 \(\mathfrak m_1,\ldots,\mathfrak m_r\)，\(M\) 为有限生成 \(R\) 模。证明
\[
\ell_R(M)=\sum_{i=1}^r\ell_{R_{\mathfrak m_i}}(M_{\mathfrak m_i}).
\]
\end{exercise}
\begin{solution}
由\cref{b3:thm:artin-decomposition}，\(R=\prod_i e_iR\)，且
\(M=\bigoplus_i e_iM\)、\(e_iM\cong M_{\mathfrak m_i}\)。
一个 \(e_iM\) 的 \(R\) 子模恰为其 \(e_iR\) 子模，因为其他坐标作用为零；因此两种底环下的合成列及长度相同。分别取各分量的合成列，再在直和中依次添入各分量，便得到长度为各长度之和的 \(M\) 合成列，证明所需等式。
\end{solution}

\begin{exercise}\label{b3:ex:03-nonnoether-subring-local}
设 \(k\) 为域，令 \(A=k+xk[x,y]\)、\(I=xk[x,y]\)。证明局部环 \(A_I\) 仍不是 Noether 环。
\end{exercise}
\begin{solution}
\(A/I=k\)，故 \(I\) 极大，\(A_I\) 的极大理想为 \(IA_I\)。由局部化正合性及乘积的分式公式，
\[
IA_I/(IA_I)^2\cong (I/I^2)_I.
\]
\(I/I^2\) 被 \(I\) 零化，且 \(A\setminus I\) 的作用只是非零常数的作用，已经可逆，所以右侧仍为以 \(xy^j+I^2\)、\(j\ge0\)，为基的无限维 \(k\) 向量空间。若 \(IA_I\) 有限生成，其模极大理想后的向量空间就有限维，矛盾。因此该局部环不是 Noether 环。
\end{solution}

\begin{exercise}\label{b3:ex:03-finite-gaussian-algebra}
证明 \(\mathbb Z[i]\) 是 Noether 环，并分解 \(\mathbb Z[i]/(5)\) 为 Artin 局部环的直积。
\end{exercise}
\begin{solution}
\(\mathbb Z[i]\) 由 \(1,i\) 作为 \(\mathbb Z\) 模生成，而整数环为 Noether 环，所以它是有限代数，从而是 Noether 环。模掉 \(5\) 后，
\[
\mathbb Z[i]/(5)\cong\mathbb F_5[T]/(T^2+1)
=\mathbb F_5[T]/\bigl((T-2)(T+2)\bigr).
\]
两个一次因子互素，由中国剩余定理得到 \(\mathbb F_5\times\mathbb F_5\)。这两个域就是局部因子，商环约化且长度为 \(2\)。
\end{solution}

\begin{exercise}\label{b3:ex:03-pointwise-noether}
设 \(k\) 为域，令 \(R=\prod_{n\ge1}k\)。对 \(f=(f_n)\in R\)，记 \(E=\{n\ge1:f_n\ne0\}\)。计算 \(R_f\)，并证明它是 Noether 环当且仅当 \(E\) 有限（允许 \(E\) 为空）。由此证明，不存在由满足 \(R_{f_i}\) 为 Noether 环的基本开集 \(D(f_i)\) 组成的有限覆盖。解释这与正文中对任意素理想 \(\mathfrak p\) 所证的 \(R_{\mathfrak p}\) 都是域为何相容。
\end{exercise}
\begin{solution}
向 \(E\) 坐标的投影把 \(f\) 送到可逆元素，逆逐坐标为 \(f_n^{-1}\)，故诱导同态
\[
R_f\longrightarrow\prod_{n\in E}k,\qquad
a/f^r\longmapsto(a_n/f_n^r)_{n\in E}.
\]
它满射，因为任意 \(E\) 坐标序列都能在其余坐标补零，成为 \(R\) 中的元素。若其值为零，\(a\) 的所有 \(E\) 坐标均为零，从而 \(fa=0\)，分式零判据给出 \(a/f^r=0\)。因此该映射是同构。\(E=\varnothing\) 时 \(f=0\)，两侧都是零环。

\(E\) 有限时，右侧是有限个域的直积，为 Artin 环，因而 Noether；\(E\) 无限时，选两两不同的坐标 \(n_1,n_2,\ldots\in E\)，由前 \(j\) 个单坐标幂等元生成的理想组成严格升链，故不是 Noether 环。

若 \(D(f_1),\ldots,D(f_s)\) 覆盖全谱，则 \((f_1,\ldots,f_s)=R\)。在每个坐标 \(n\) 上，至少一个 \(f_i\) 的值必须非零，否则这几个元素的任何线性组合在该坐标都为零，不可能等于 \(1\)。因此各 \(E_i\) 的并必须是全部正整数，有限个有限集不能满足此条件。于是 Noether 的基本开局部化不能给出这样的有限覆盖。若每个点都有这样的 Noether 基本开邻域，全谱的拟紧性就会给出一个有限子覆盖，与上面的结论矛盾。正文的域结论对每个素理想都成立，却不能保证各点有 Noether 的基本开邻域；这正是\cref{b3:prop:finite-presentation-open-cover}所需的额外条件。
\end{solution}

\begin{exercise}\label{b3:ex:03-noether-artin-distinction}
分别给出下列例子并验证：Noether 而非 Artin 的环；Artin 而非 Noether 的模；有限生成但非有限展示的模。最后说明为什么不存在交换 Artin 而非 Noether 的环。
\end{exercise}
\begin{solution}
\(\mathbb Z\) 的每个理想为主理想，故是 Noether 环；理想链
\(2^n\mathbb Z\) 严格下降，故不是 Artin 环。
取素数 \(p\)，整数模 \(\mathbb Z[1/p]/\mathbb Z\) 的子模分类已在\subsecref{b3:subsec:0304:03}证明：每个真子模有限循环，所以降链稳定；但它含有阶依次为 \(p^n\) 的严格上升子群链，所以不是 Noether 模。
取域 \(k\)，令 \(R=k[x_1,x_2,\ldots]\)、\(I=(x_1,x_2,\ldots)\)，则 \(R/I\) 循环而有限生成；若它有限展示，\cref{b3:prop:finite-presentation-kernel}会迫使 \(I\) 有限生成，与变量数的论证矛盾。
最后一种环由\cref{b3:thm:artin-finite-length}排除：交换 Artin 环 \(R\) 作为自身的模具有有限长度，必满足升链条件。
\end{solution}
